\section{问题二模型、求解与结果}

\subsection{无泄漏预测与联合场景}

主预测采用目标日前28个完整历史日的加权平均，时间半衰期为7日，同星期权重乘2。对每个历史日保存负荷与光伏的整日样本外残差对，不独立打乱144个时段，也不拆开两类变量。目标日场景由最近最多28个成对残差日等概率枚举：
\begin{equation}
 L_{i,t}^{\omega}=\max\{0,\widehat L_{i,t}+\varepsilon_{j(\omega),t}^{L}\},\quad
 G_{i,t}^{\omega}=\max\{0,\widehat G_{i,t}+\varepsilon_{j(\omega),t}^{G}\}.
\end{equation}
该构造借鉴概率轨迹与滚动回放的思想\cite{lemos2022}，保留日内连续性及负荷与光伏的联合偏差。2月1日共有24个可用残差场景，随后最多使用28个；所有预测与场景只由目标日以前的数据产生。

\subsection{两阶段随机规划}

第一阶段在0:00确定$\bm x_i=(\bm q_i,\bm c_i,\bm d_i,\bm s_i,\bm z_i)$，并在全部场景间共享。第二阶段仅包含场景相关的紧急购电$e_{i,t}^{\omega}$与富余量$u_{i,t}^{\omega}$。日模型为
\begin{equation}
\min_{\bm x_i}\ \sum_t p_tq_{i,t}
+\sum_{\omega\in\Omega_i}\pi_\omega\sum_t5p_te_{i,t}^{\omega}
+\Phi_i(s_{i,145}),
\label{eq:q2obj}
\end{equation}
并对每个场景满足式\eqref{eq:balance}。终端价值取$\Phi_i(s)=-vs$，其中$v=\min_t(p_t)/\eta_c=0.412555556$元/电池内部kWh，仅用于近似跨日库存价值，不计入实际账单。1月采用电池闲置热身，2月1日储能量为6000 kWh；2月至12月逐日连续回放，12月31日末约束为6000 kWh。

取得计划后，用当天实际曲线计算
\begin{equation}
 e_{i,t}^{\mathrm{real}}=\pos{L_{i,t}^{\mathrm{real}}+c_{i,t}-q_{i,t}-G_{i,t}^{\mathrm{real}}-d_{i,t}},
\end{equation}
实际费用为$F_i^{\mathrm{real}}=\sum_tp_tq_{i,t}+\sum_t5p_te_{i,t}^{\mathrm{real}}$。因此全年过程是滚动求解334个日内两阶段随机MILP，并非334阶段随机规划。

\subsection{全年结果与对照}

\begin{table}[H]
\centering
\caption{问题二334天方法对照}
\label{tab:q2-comparison}
\begin{tabularx}{\textwidth}{Xrrr}
\toprule
方法 & 实际总费/万元 & 紧急购电费/万元 & 紧急购电量/MWh\\
\midrule
加权历史+场景规划 & 1676.65 & 121.67 & 320.32\\
7日均值+场景规划 & 1679.73 & 112.13 & 296.68\\
加权历史+单一预测 & 2405.78 & 1237.18 & 2995.82\\
同场景+无储能 & 2054.28 & 108.40 & 264.34\\
\bottomrule
\end{tabularx}
\end{table}

主方案相对相同加权预测的单一场景规划费用降低30.308\%，说明在5倍紧急电价下显式处理预测误差具有显著经济作用；相对同场景无储能方案降低18.383\%，表明储能跨时段调节有效。主方案相对7日均值场景方案仅低0.184\%，且其紧急购电量更高，不能将这一小幅费用差异解释为预测精度全面占优。

两种场景方案的负荷MAE分别为739.966 kW和781.211 kW，光伏MAE分别为177.805 kW和153.400 kW；加权历史对负荷更准，7日均值对光伏更准。这进一步说明预测误差与调度费用必须同时评价。

\begin{figure}[H]
  \centering
  \includegraphics[width=\textwidth]{q2_comparison.png}
  \caption{问题二实际费用构成及紧急购电量对照。所有方案使用相同334天和初末储能条件。}
  \label{fig:q2-comparison}
\end{figure}

四个指定日期中，3月20日、6月21日、9月23日和12月21日的实际费用分别为44129.72、50918.55、49971.63和65868.92元，对应紧急购电量分别为252.83、0、2072.26和762.57 kWh。9月23日应急需求显著高于其余三日，说明全年总量之外仍需检查季节与日内短缺分布。

\begin{figure}[H]
  \centering
  \includegraphics[width=\textwidth]{q2_selected_days.png}
  \caption{问题二四个指定日期的计划购电与实际紧急购电轨迹}
  \label{fig:q2-days}
\end{figure}

\subsection{敏感性与计算复核}

终端价值系数取基准的0.8、1.0、1.2倍时，全年实际费用变化分别为$0.000\%$、$0$、$-0.0527\%$；历史残差幅度取0.8、1.0、1.2倍时，费用变化分别为$+1.7847\%$、$0$、$+1.3917\%$。每个点均重新连续回放334天，该局部试验只支持“基准附近有一定稳定性”，不能推出全局稳健。

程序完成20项单元检查和8组全年独立复核，包括日期与时间索引、能量守恒、SOC边界、充放电互斥、费用重算、模板读回及防未来信息泄漏。主结果采用整体MILP；Benders方法在给定20 s和200轮限制内未取得所需收敛证据，故未用于全年结论。

需要指出，主方案有333天日末储能量贴近1200 kWh下限，表明当前线性终端价值较弱。该现象不破坏物理可行性，但限制了跨日运行解释；后续应比较48 h滚动时域或更精细的库存价值函数。
